\scanpage{4}\mspage{2}

sur $\unsure{E|\mathcal E}^{(\unsure{\nu|v})}$ une structure d'\deleted{algèbre}anneau gradué ?

Il suffit que l'on ait les relations (1.1.)
\[
 \deleted{\unsure{*}}\left\{\begin{aligned}
 \pi_i^2&=\pi_i\\
 \pi_i\pi_j&=\pi_j\pi_i=0\quad\text{si }i\ne j\\
 \sum\pi_i&=1
 \end{aligned}\right.
\]
qui impliquent que les $\operatorname{pr}_{\nu}^{E}$ définis par (1.2.) sont des projections deux à deux orthogonales de somme 1, donc que les $E^{(\nu)}$ définissent bien une graduation \inserted{de $k$-module} sur $E^{(\nu)}$, \deleted{\unsure{puis}} \inserted{et} que $E^{(\nu)} E^{(\nu')}\subset E^{(\nu+\nu')}$. Noter aussi que (1.4.) que (1.4) reste valable, car \deleted{\unsure{comme on a pour tt} $u$} \inserted{on aura la décomposition}
% slanted
\margin{on peut remplacer $\mathbb Z$ par un \uline{groupe} quelconque}
\[
 \deleted{u=\sum_{ij}\pi_j u\pi_i}
\]
\deleted{\unsure{d'où la relation|de la relation}}
\[
 \deleted{\pi_{i+\nu}u\pi_i=\pi_{i+\nu}u\pi_i}\qquad E_i^j
\]
\[
 \left\{\begin{aligned}
 E&=\bigoplus_{i,j}\pi_j E\pi_i\\
 \operatorname{pr}_{i,j}(u)&=\pi_j u\pi_i
 \end{aligned}\right.\qquad\text{donc }\ u=\sum_{\alpha,\beta}\pi_\beta u\pi_\alpha
\]
et la relation $u\pi_i=\pi_{i+\nu}u$ implique, puisque $u\pi_i=\sum_{\beta}\pi_{\beta}u\pi_i$, et $\pi_{\unsure{j|i}+\nu}u=\sum_i\pi_{j+\nu}u\pi_\alpha$, que $\pi_\beta u\pi_i=0$ si $\beta\ne i+\nu$.

\margin{Regarder les modules sur \deleted{\unsure{un|des}} anneau $E$ muni \deleted{\unsure{de}} \deleted{\unsure{des}} d'éléments $\pi_i$ comme dessus.}

3. \deleted{\unsure{Nous notons|Nous pouvons}} Partons enfin d'une \deleted{algèbre} \inserted{anneau} graduée $E^*$ \deleted{sur $k$}, et posons nous la question de l'existence et \uline{l'unicité} d'une famille $(\pi_i)_{i\in\mathbb Z}$ d'élts de $E^*$ \inserted{satisfaisant (1.1.)}, définissant la graduation par (1.3.) (\deleted{\unsure{donc}} \inserted{ou} par (1.4)). On \deleted{\unsure{supp}} supposera ici que
% archive sheet number: 3
